

\begin{frame}[fragile]{De SQL vers une forme opératoire: l'algèbre}

Titre des films parus en 1958, où l'un des acteurs joue
le rôle de John Ferguson.

\vfill

En SQL:

 \begin{minted}[frame=leftline,
               baselinestretch=.9,
               fontsize=\small,
               framesep=2mm]{sql}
    select titre
    from   Film f, Role r
    where  nom_role ='Ferguson'
    and    f.id = r.id_ilm
    and    f.annee = 1958
\end{minted}

\vfill

\begin{tabular}{p{9cm}p{5cm}}
Deux sélections (l'année, le rôle), une jointure, une projection.
&
~
\end{tabular}


\end{frame}


\begin{frame}[fragile]{Le Plan d'Exécution Logique (PEL)}


\only<1-1|handout:1-1>{
En algèbre, sous une forme normalisée projection-sélection-jointure:
\[
  \pi_{titre}(\sigma_{annee=1958 \land nom\_role='\mathrm{Ferguson}'}(Film \Join_{id=id\_film} Role))
\]
}

\only<2-3|handout:2-3>{
La même expression, sous forme d'un arbre.

\vfill
\begin{small}
\psset{levelsep=1.2cm,treesep=0.5cm}
\IteratorSansNom{$\pi_{titre}$}{}{
  \IteratorSansNom{$\sigma_{annee=1958 \land nom\_role='Ferguson'}$}{}{
     \IteratorSansNom{$\Join$}{}{ 
       \Database{Film}{}
       \Database{Role}{}
    }
  }
}
\end{small}
}
\only<3-3|handout:3-3>{
\vfill

\begin{tabular}{p{9cm}p{5cm}}
Est-ce le seul? Non: des \alert{réécritures} vont nous permettre d'en trouver
d'autres.
&
~
\end{tabular}
}
\end{frame}



\begin{frame}{Réécritures}

Il y a plusieurs expressions
\alert{équivalentes} pour une même requête. 

\smallskip

On peut explorer ces expressions grâce à des règles de réécriture.
Exemple


\begin{redlist}
\item \alert{Commutativité des jointures}\,:        $ R \Join S  \equiv S \Join R$

\item \alert{Regroupement des sélections}\,: 
        $\sigma_{A='a' \wedge B='b'}(R) \equiv
        \sigma_{A='a'}(\sigma_{B='b'}(R))$
        
\item \alert{Commutativité de $\sigma$ et de $\pi$}\,: 
        $\pi_{A_1, A_2, \ldots A_p}(\sigma_{A_i='a'} (R)) \equiv
        \sigma_{A_i='a'}(\pi_{A_1, A_2, \ldots A_p}(R))$

        \item \alert{Commutativité de $\sigma$ et de $\Join$}\,: 
$\sigma_{A='a'} (R[\ldots A\ldots] \Join S) \equiv
     \sigma_{A='a'}(R) \Join S$
\item etc.
\end{redlist}

\vfill



\begin{tabular}{p{9cm}p{5cm}}
L'application dirigée d'une règle d'équivalence (réécriture)
transforme une expression $e$ 
en une expression $e'$ \alert{équivalente}.
&
~
\end{tabular}

\end{frame}


\begin{frame}{Ce que fait l'optimiseur}

Trouve les expressions équivalentes, 
évalue leur coût et choisit la meilleure.

\vfill
\alert{Important}:
on ne peut pas \alert{énumérer} tous les plans possibles (trop long):
on applique des \alert{heuristiques}. 

\vfill

\begin{tabular}{p{9cm}p{5cm}}
Heuristique classique: \alert{réduire la taille
des données}
\begin{redbullet}
\item en filtrant les nuplets par
 des sélections
 \item en les simplifiant par des projections
\end{redbullet}
\alert{dès que possible.}
&
~
\end{tabular}

\end{frame}



%%%%%%%%%%%%%%%%%%%%%%%%%%%%%%%%%%%%%%%%%%%%%%%%%%%%%%%

\begin{frame}{Illustration}

\only<1-1|handout:1-1>{
Reprenons: le film
paru en 1958 avec un rôle 'Ferguson'.

\vfill
\begin{small}
\psset{levelsep=1cm,treesep=0.5cm}
\IteratorSansNom{$\pi_{titre}$}{}{
  \IteratorSansNom{$\sigma_{annee=1958 \land nom\_role='Ferguson'}$}{}{
     \IteratorSansNom{$\Join$}{}{ 
       \Database{Film}{}
       \Database{Role}{}
    }
  }
}
\end{small}
}


\only<2-2|handout:2-2>{
\bigskip
Première étape: je sépare les sélections.

\vfill

\begin{small}
\psset{levelsep=1cm,treesep=0.5cm}
\IteratorSansNom{$\pi_{titre}$}{}{
  \IteratorSansNom{$\sigma_{annee=1958}$}{}{
    \IteratorSansNom{$\sigma_{nom\_role='Ferguson'}$}{}{
     \IteratorSansNom{$\Join$}{}{ 
       \Database{Film}{}
       \Database{Role}{}
    }
   }
  }
}
\end{small}
}
\only<3-|handout:3->{
Puis on pousse chaque sélection vers sa table.
\vfill
\begin{small}
\psset{levelsep=1cm,treesep=0.5cm}
\IteratorSansNom{$\pi_{titre}$}{}{
     \IteratorSansNom{$\Join$}{}{ 
       \IteratorSansNom{$\sigma_{annee=1958}$}{}{
       \Database{Film}{}
       }
  \IteratorSansNom{$\sigma_{nom\_role='Ferguson'}$}{}{
       \Database{Role}{}
    }
   }
}
\end{small}
}
\only<4-|handout:4->{
\vfill
\begin{tabular}{p{9cm}p{5cm}}
À ce stade: reste à choisir les chemins d'accès et les algorithmes de jointure.

&
~
\end{tabular}
}
\end{frame}

%%%%%%%%%%%%%%%%%%%%%%%%%%%%%%%%%%%%%%%%%%%%%%%%%%%%%%%
\begin{frame}{Résumé: la réécriture algébrique}

Principes essentiels\,:

\begin{redlist}
  \item L'algèbre permet d'obtenir une version opératoire de la requête.
      
   \item Les équivalences algébriques permettent d'explorer
    un ensemble de plans.
      
    \item L'optimiseur évalue le coût de chaque plan.

\end{redlist}

\vfill

Bien retenir

\begin{tabular}{p{9cm}p{5cm}}
\begin{redbullet}
  \item \alert{Heuristique}: on ne peut pas \alert{tout explorer}
  \item \alert{Nécessaire mais pas suffisant}: il reste à choisir 
   le bon algorithme pour chaque opération.
\end{redbullet}
&
~
\end{tabular}


%MUn opérateur bloquant a de forts besoins  en mémoire, et un temps de réponse
%qui peuvent être importants.

\only<2|handout:2>{
\vfill
\centerline{\alert{Merci!}}
}


\end{frame}
